\documentclass[11pt,a4paper]{article}
\usepackage{fontspec}
\usepackage{xeCJK}
\setCJKmainfont{STSong}
\setCJKsansfont{STHeiti}
\usepackage{amsmath,amssymb,amsthm}
\usepackage{mathtools}
\usepackage{bm}
\usepackage{booktabs}
\usepackage{array}
\usepackage{enumitem}
\usepackage{geometry}
\usepackage{xcolor}
\usepackage{tcolorbox}
\usepackage{algorithm}
\usepackage{algpseudocode}
\usepackage{pifont}
\usepackage[pdfencoding=auto,psdextra]{hyperref}

\geometry{margin=2.5cm}

\newcommand{\loss}{\mathcal{L}}
\newcommand{\E}{\mathbb{E}}
\newcommand{\R}{\mathbb{R}}
\newcommand{\KL}{\mathrm{KL}}
\newcommand{\N}{\mathcal{N}}
\newcommand{\sg}{\mathrm{sg}}
\newcommand{\Enc}{\mathcal{E}}
\newcommand{\Dec}{\mathcal{D}}
\newcommand{\Xspace}{\mathcal{X}}
\newcommand{\Zt}{\mathcal{Z}_T}
\newcommand{\Zs}{\mathcal{Z}_S}
\newcommand{\Tmap}{\mathsf{T}}
\newcommand{\push}{\#}
\newcommand{\vT}{v_{T}}
\newcommand{\vS}{v_{\theta}}
\newcommand{\Jac}{\mathbf{J}}
\newcommand{\norm}[1]{\left\lVert #1 \right\rVert}

\theoremstyle{definition}
\newtheorem{definition}{Definition}
\newtheorem{remark}{Remark}
\newtheorem{proposition}{Proposition}
\newtheorem{corollary}{Corollary}
\newtheorem{lemma}{Lemma}
\newtheorem{theorem}{Theorem}
\newtheorem{assumption}{Assumption}

\title{\textbf{跨 VAE 隐空间的 Cross-Model On-Policy Distillation}\\[0.3em]
\large XOPD with Heterogeneous Latent Spaces：理论形式化与算法族\\[0.4em]
\large 以 FLUX.2-dev (32B) $\to$ SD3.5 为参考实例}
\author{Flow Factory}
\date{}

\begin{document}
\maketitle

\begin{abstract}
现有 XOPD 蒸馏器（\texttt{trainers/xopd/trainer.py}）建立在一个硬约束之上：
teacher 与 student \emph{共享同一个 VAE 隐空间}（\texttt{load\_teacher\_transformer}
显式断言 \texttt{teacher.in\_channels == student.in\_channels}）。当我们希望把
FLUX.2-dev（其 VAE 隐空间 $\Zt$）蒸馏到 Stable Diffusion 3.5（隐空间 $\Zs$，
通道数/分辨率/统计均不同）时，这个约束被打破：teacher 的 velocity 场与 transition
分布都活在 $\Zt$ 中，而 student 的 rollout、loss、梯度都活在 $\Zs$ 中。

本文把"跨 VAE 蒸馏"形式化为一个\textbf{隐空间传输映射} (latent transport map)
$\Tmap:\Zt\to\Zs$ 的设计问题，并给出一个统一框架：XOPD 的两个 loss（L0
velocity regression 与 L1 transition-KL）分别对 $\Tmap$ 提出\emph{不同强度}的要求
——L0 只需在\emph{干净样本}上传输（sample transport），而 L1 需要传输整个
\emph{velocity 场}（vector-field pushforward），后者只有当 $\Tmap$ 是可逆光滑映射、
且其 Jacobian 行为良好时才成立。由此我们得到一个关键结论：\textbf{仿射/线性传输是
唯一能让 L0 与 L1 同时精确成立的传输类}，因为它的 Jacobian 是常数、且保持 flow
matching 的插值路径结构。

在此框架下我们形式化并比较四类算法：(M1) pixel-space 桥接（decode-encode），
(M2) 学习式 latent projector（线性 / 仿射 / 非线性），(M3) CV-VAE 式重建对齐
(\href{https://arxiv.org/abs/2405.20279}{arXiv:2405.20279}，把对齐推到 VAE 训练阶段)，
以及我们新提出的 (M4) 概率耦合 / 分布矩匹配、(M5) cycle-consistent 联合学习传输、
(M6) shared-decoder 微调。对每种方法我们给出其对 L0/L1 的精确作用、有损性来源、
可微性与梯度路径、训练开销，并用一个统一的\emph{传输一致性误差}
$\delta_{\Tmap}$ 给出 loss 偏差的上界。最后给出与 \texttt{trainers/xopd} 现有代码
对齐的实现设计（抽象出 \texttt{VAETransport}、放松 \texttt{in\_channels} 断言、
传输在 L0/L1 中的插入点、config 接口与实现优先级），\textbf{本文不写生产代码，只做
设计}。
\end{abstract}

\tableofcontents
\newpage

%% ============================================================
\section{Motivation 与问题设定}
\label{sec:motivation}
%% ============================================================

\subsection{XOPD 回顾：共享隐空间假设}

设 student 是一个在隐空间 $\Zs\subseteq\R^{d_S}$ 上的 flow-matching 模型，其
velocity 场为 $\vS(z,t)$（$\theta$ 为可训练参数，实现中是 LoRA）。teacher 是一个
冻结的大模型，velocity 场为 $\vT$。XOPD 的两个 loss（见
\texttt{cross\_opd\_xopd.md}）为：

\begin{itemize}[leftmargin=1.4em]
\item \textbf{L0（velocity regression，off-policy）}：teacher 先 rollout 出干净
隐变量 $z_0^{T}$，构造数据路径 $z_t=(1-\sigma_t)z_0^T+\sigma_t\,\epsilon$，然后在
\emph{同一个} $z_t$ 上回归
\begin{equation}
\loss_{\mathrm{L0}}
=\E_{t,\epsilon}\Big[\,w(t)\,\big\lVert \vS(z_t,t)-\vT(z_t,t)\big\rVert^2\,\Big].
\label{eq:l0-shared}
\end{equation}

\item \textbf{L1（transition matching，on-policy）}：student 自己 rollout 轨迹
$\{x_k\}$，对每个训练时间步计算 student / teacher 的转移高斯均值
$\mu_\theta(x_k),\mu_T(x_k)$，以闭式高斯 KL（即 $D_k$）作为 pathwise loss：
\begin{equation}
D_k=\frac{\lVert \mu_\theta(x_k)-\mu_T(x_k)\rVert^2}{2\bar\sigma_k^2},
\qquad
\loss_{\mathrm{L1}}=\E\Big[\textstyle\sum_k D_k\Big]
\;(+\,\text{REINFORCE}\,+\,\text{KL anchor}).
\label{eq:l1-shared}
\end{equation}
\end{itemize}

\eqref{eq:l0-shared} 与 \eqref{eq:l1-shared} 的共同前提是 $\vS$ 与 $\vT$
\textbf{定义在同一个隐空间}：L0 在同一个 $z_t$ 上比较两者，L1 在 student 访问到的
状态 $x_k\in\Zs$ 上查询 teacher 的均值 $\mu_T(x_k)$。代码用一行断言保证它：

\begin{tcolorbox}[colback=red!4,colframe=red!50!black,title=\texttt{flux2\_klein.py:1363} 的硬约束]
\small
\begin{verbatim}
student_in_channels = self.pipeline.transformer.config.in_channels
teacher_in_channels = teacher.config.in_channels
if teacher_in_channels != student_in_channels:
    raise ValueError("... shared-latent-space assumption is violated ...")
\end{verbatim}
\end{tcolorbox}

\subsection{跨 VAE 问题}

记 teacher VAE 为 $(\Enc_T,\Dec_T)$、隐空间 $\Zt\subseteq\R^{d_T}$；student VAE 为
$(\Enc_S,\Dec_S)$、隐空间 $\Zs\subseteq\R^{d_S}$；共享的像素空间为 $\Xspace$。
$\Enc_T:\Xspace\to\Zt$，$\Dec_T:\Zt\to\Xspace$，student 同理。一般地
\[
d_T\neq d_S,\quad
\Zt\not\cong\Zs,\quad
\Enc_T\neq\Enc_S,
\]
通道数、空间下采样率、逐通道统计（FLUX.2 VAE 还带一个 BatchNorm 标准化层，见
\texttt{decode\_latents}）都不同。直接把 $\vT$ 代入 \eqref{eq:l0-shared} 或把
$\mu_T(x_k)$ 用于 \eqref{eq:l1-shared} 在类型上就不合法。

\begin{tcolorbox}[colback=blue!4,colframe=blue!50!black,title=核心问题]
我们需要一个\textbf{隐空间传输映射} $\Tmap:\Zt\to\Zs$，把 teacher 在 $\Zt$ 中产生
的信号（干净样本 $z_0^T$、或整个 velocity 场 $\vT$）搬运到 $\Zs$ 中，使得 XOPD 的
监督信号重新变得"类型合法且语义正确"。本文的全部内容就是：$\Tmap$ 应满足什么性质、
有哪些构造法、各自对 L0/L1 的作用与代价。
\end{tcolorbox}

%% ============================================================
\section{形式化：隐空间传输映射}
\label{sec:formal}
%% ============================================================

\subsection{两种传输：样本传输 vs.\ 速度场传输}

跨 VAE 蒸馏的微妙之处在于：L0 与 L1 需要的\emph{不是同一种}传输。

\begin{definition}[语义一致性 / 像素锚定]
我们称 $\Tmap:\Zt\to\Zs$ 是\textbf{像素一致}的，若对 teacher 隐变量 $z\in\Zt$，
$\Tmap(z)$ 解码出的图像与 $z$ 解码出的图像一致：
\begin{equation}
\Dec_S\big(\Tmap(z)\big)\approx \Dec_T(z),
\qquad\forall z\in\Zt.
\label{eq:pixel-consistent}
\end{equation}
像素一致是所有合理传输的\emph{语义正确性}判据：两个隐变量"代表同一张图"。
\end{definition}

\begin{definition}[样本传输与速度场传输]
\leavevmode
\begin{itemize}[leftmargin=1.4em]
\item \textbf{样本传输 (sample transport)}：只要求把\emph{干净样本} $z_0^T$ 搬到
$\Zs$：$z_0^S=\Tmap(z_0^T)$。L0 只需要这个。
\item \textbf{速度场传输 (vector-field pushforward)}：要求把整个 velocity 场
$\vT$ 搬到 $\Zs$ 上的一个场 $(\Tmap_\push \vT)$，使其在 $\Zs$ 中生成的概率流与
teacher 在 $\Zt$ 中的流经 $\Tmap$ 推前后一致。L1 需要这个。
\end{itemize}
\end{definition}

\begin{remark}[隐变量\emph{不}服从高斯分布——四个分布的辨析]
\label{rem:not-gaussian}
本文从不假设 latent 服从高斯分布，这一点必须澄清，因为 VAE 语境里至少有四个分布被
经常混淆，它们的高斯性截然不同：
\begin{itemize}[leftmargin=1.5em]
\item \textbf{逐图后验} $q(z\mid x)=\N(\mu_\phi(x),\mathrm{diag}\,\sigma_\phi^2(x))$：
\emph{高斯}（reparameterization trick 构造使然，单张图）；
\item \textbf{聚合后验/边际} $q(z)=\int q(z\mid x)\,p_{\mathrm{data}}(x)\,dx$：
\emph{非高斯}（一张图一个分量的巨型混合，多峰、携带全部图像结构）；
\item \textbf{干净隐变量} $z_0$ 的分布：\emph{非高斯}（$=$ 上面的边际，正是 flow 要去
刻画的目标——若它高斯则生成平凡）；
\item \textbf{噪声端} $z_1\sim\N(0,I)$：\emph{高斯}（flow matching 的 prior 端，定义）。
\end{itemize}
所以"干净图像/干净 latent 不服从高斯"是\textbf{正确}的：非高斯的是\emph{边际}
$q(z_0)$。唯一天然高斯的是\emph{逐图后验}与\emph{噪声端}。此外，要让\emph{边际}接近
$\N(0,I)$ 还需 KL 正则权重足够大且无 prior hole；而 FLUX.2 / SD3.5 这类扩散用 VAE 的
KL 权重极小（$\sim 10^{-6}$，近乎纯自编码器以保重建保真度），其边际\textbf{比教科书
VAE 离高斯更远}。FLUX.2 VAE 的 BatchNorm（\texttt{decode\_latents} 中
$z\cdot\mathrm{std}+\mathrm{mean}$）只标准化\emph{逐通道一阶/二阶矩}，标准化 $\neq$
高斯化——通道内形状、跨通道/跨空间依赖仍非高斯。

\textbf{对本文的影响}：偏好线性传输的真正依据是\emph{几何}（Prop.~\ref{prop:affine}，
见其后的 Rem.~\ref{rem:affine-geometry}），\textbf{不}依赖任何分布假设；L1 的逐步高斯
KL（$D_k$）匹配的是\emph{条件}转移分布 $z_{k+1}\mid z_k$（由 SDE 加性噪声构造，本就
高斯），与 $z_0$ 边际是否高斯无关。下面 M4（\S\ref{sec:methods}）只在边际\emph{恰好}
近高斯时才升级为全局最优；一般情形它退化为\emph{矩匹配}初始化。
\end{remark}

\subsection{速度场的推前 (pushforward)}

\begin{proposition}[ODE 流的推前]
\label{prop:pushforward}
设 $\Tmap:\Zt\to\Zs$ 为 $C^1$ 微分同胚，Jacobian 记为
$\Jac_{\Tmap}(z)=\partial \Tmap/\partial z$。若 $z^T(t)$ 满足
teacher 的概率流 ODE $\dot z^T=\vT(z^T,t)$，令 $z^S(t)=\Tmap\big(z^T(t)\big)$，则
$z^S(t)$ 满足
\begin{equation}
\dot z^S
=\underbrace{\Jac_{\Tmap}\!\big(\Tmap^{-1}(z^S)\big)\;
\vT\big(\Tmap^{-1}(z^S),t\big)}_{\displaystyle =:(\Tmap_\push\vT)(z^S,t)}.
\label{eq:pushforward}
\end{equation}
即推前速度场 $\Tmap_\push\vT$ 由链式法则 $\dot z^S=\Jac_{\Tmap}\dot z^T$ 给出。
\end{proposition}

\begin{proof}
链式法则 $\dot z^S=\Jac_{\Tmap}(z^T)\dot z^T=\Jac_{\Tmap}(z^T)\vT(z^T,t)$，再代入
$z^T=\Tmap^{-1}(z^S)$。
\end{proof}

\begin{corollary}[L1 需要 Jacobian 与可逆性]
\label{cor:l1-hard}
L1 在 student 状态 $x_k\in\Zs$ 上需要 teacher 的"意见"
$(\Tmap_\push\vT)(x_k,t)$。由 \eqref{eq:pushforward}，这要求：(i) $\Tmap$ \emph{可逆}
（取 $\Tmap^{-1}(x_k)$ 才能查询 teacher），(ii) $\Tmap$ \emph{可微}且 Jacobian 可得。
对一般非线性 $\Tmap$ 这两点都昂贵；对像素桥接 $\Tmap=\Enc_S\circ\Dec_T$，$\Tmap^{-1}$
无解析式、Jacobian 需 VJP/JVP，逐步代价极高。
\end{corollary}

\subsection{为什么线性 / 仿射传输是"甜区"}

\begin{proposition}[仿射传输保持 flow-matching 路径与速度]
\label{prop:affine}
设 $\Tmap(z)=Az+b$（$A\in\R^{d_S\times d_T}$）。则：
\begin{enumerate}[leftmargin=1.6em,label=(\alph*)]
\item \textbf{保插值路径}：teacher 的 flow-matching 路径
$z_t^T=(1-\sigma_t)z_0^T+\sigma_t\,\epsilon^T$ 经 $\Tmap$ 推前后仍是 student 空间里
的一条 flow-matching 路径
$z_t^S=(1-\sigma_t)\,\Tmap(z_0^T)+\sigma_t\,(A\epsilon^T)$，端点为
$\Tmap(z_0^T)$ 与噪声 $A\epsilon^T$；
\item \textbf{速度常数倍传输}：Jacobian $\Jac_{\Tmap}\equiv A$ 为常数，故
$(\Tmap_\push\vT)(z^S,t)=A\,\vT(A^{+}(z^S-b),t)$（$A^{+}$ 为伪逆，$A$ 行满秩可逆时为
逆），速度场的传输\emph{无需}逐点 Jacobian 计算。
\end{enumerate}
\end{proposition}

\begin{proof}
(a) 直接代入：$\Tmap(z_t^T)=A\big((1-\sigma_t)z_0^T+\sigma_t\epsilon^T\big)+b
=(1-\sigma_t)(Az_0^T+b)+\sigma_t(A\epsilon^T)+\big[(1-\sigma_t)+\sigma_t-1\big]b$，
末项为 $0$。(b) 由 Prop.~\ref{prop:pushforward} 及 $\Jac_{\Tmap}\equiv A$。
\end{proof}

\begin{remark}[偏好线性传输是\emph{几何}结论，不依赖分布假设]
\label{rem:affine-geometry}
Prop.~\ref{prop:affine} 的两条性质对\textbf{任意}干净隐变量 $z_0^T$ 成立，证明只用了
仿射映射的代数性质（$\Tmap$ 与凸组合可交换、Jacobian 为常数），\emph{完全没有}用到
$z_0^T$ 服从高斯或任何特定分布（呼应 Rem.~\ref{rem:not-gaussian}）。这才是本文偏好
线性传输的\textbf{真正理由}：它在计算上干净（速度场推前 $=A\vT$，无需逐点 Jacobian；
插值路径保结构），与 latent 边际是否高斯无关。是否\emph{够用}（线性 $\delta_{\Tmap}$
是否够小）是关于"同一批图训出的两个 VAE 隐空间是否近似线性可通"的\emph{经验}事实
（参见 relative-representation 一类工作的经验观察），由 \S\ref{sec:formal} 的
$\delta_{\Tmap}$ 实测判定，而非由高斯定理保证。
\end{remark}

\begin{remark}[非线性传输破坏路径结构]
\label{rem:nonlinear-break}
若 $\Tmap$ 非线性，则 $\Tmap\big((1-\sigma_t)z_0^T+\sigma_t\epsilon^T\big)\neq
(1-\sigma_t)\Tmap(z_0^T)+\sigma_t\Tmap(\epsilon^T)$，即 teacher 的线性插值路径\emph{不}映成
student 空间里的线性插值路径。于是 L0 中"在 $z_t$ 上比较速度"这一构造对非线性
$\Tmap$ 失去了 flow-matching 的解析靶子；只能退化为样本传输 L0
（见 \S\ref{sec:l0l1}）。这与 CV-VAE 选择\emph{冻结 VAE、用重建 loss 把两个隐空间
"对齐成近似可线性互通"} 的动机一致。
\end{remark}

\subsection{传输一致性误差与 loss 偏差上界}

\begin{definition}[传输一致性误差]
\label{def:delta}
定义传输在像素空间的往返重建误差
\begin{equation}
\delta_{\Tmap}
=\E_{z\sim p_T}\big\lVert \Dec_S(\Tmap(z))-\Dec_T(z)\big\rVert,
\label{eq:delta}
\end{equation}
其中 $p_T$ 是 teacher 在相关时刻的隐变量分布。$\delta_{\Tmap}=0$ 即严格像素一致
\eqref{eq:pixel-consistent}。
\end{definition}

\begin{proposition}[L0 样本传输的偏差被 $\delta_{\Tmap}$ 控制]
\label{prop:bound}
设 student VAE 编码器 $\Enc_S$ 在像素空间 $L_E$-Lipschitz。用样本传输构造 L0 的
干净靶子 $z_0^S=\Tmap(z_0^T)$，并与"理想靶子" $\bar z_0^S=\Enc_S(\Dec_T(z_0^T))$
（即真正像素一致的 student 隐变量）相比，则
\begin{equation}
\E\big\lVert z_0^S-\bar z_0^S\big\rVert
\le L_E\,\delta_{\Tmap}
+\E\big\lVert \Enc_S(\Dec_T(z_0^T))-\Enc_S(\Dec_S(\Tmap(z_0^T)))\big\rVert,
\end{equation}
其中第二项在 $\Tmap$ 像素一致（\eqref{eq:pixel-consistent} 成立）时为 $0$。
因此 L0 监督靶子的偏差由传输的像素往返误差 $\delta_{\Tmap}$ 线性控制。
\end{proposition}

\begin{remark}
Prop.~\ref{prop:bound} 给出了一个干净的工程判据：\textbf{任何传输方法的"有损性"都可
以用 $\delta_{\Tmap}$（一个可离线测量的标量）来量化与比较}。这也是
\S\ref{sec:methods} 中各方法横向对比的统一坐标。
\end{remark}

%% ============================================================
\section{传输对 L0 与 L1 的作用（统一分析）}
\label{sec:l0l1}
%% ============================================================

\subsection{L0：样本传输即可，且优先于速度传输}

\begin{tcolorbox}[colback=green!4,colframe=green!50!black,title=关键洞察 1：L0 用样本传输绕开 teacher 速度]
跨 VAE 下，L0 的最自然形式\textbf{不再回归 teacher 速度}，而是：
\begin{enumerate}[leftmargin=1.5em]
\item teacher 在 $\Zt$ 中 rollout 干净样本 $z_0^T$；
\item 样本传输 $z_0^S=\Tmap(z_0^T)$（理想情形 $\Tmap=\Enc_S\circ\Dec_T$）；
\item 在 student 空间做\emph{标准 flow matching}：$z_t^S=(1-\sigma_t)z_0^S+\sigma_t\epsilon$，
回归到\emph{解析靶子}
\begin{equation}
\loss_{\mathrm{L0}}^{\times}
=\E_{t,\epsilon}\Big[w(t)\big\lVert \vS(z_t^S,t)-(\epsilon-z_0^S)\big\rVert^2\Big].
\label{eq:l0-cross}
\end{equation}
\end{enumerate}
这只需要\emph{样本传输}（Def.~\ref{def:delta} 的 $\delta_{\Tmap}$ 控制其偏差，
Prop.~\ref{prop:bound}），\textbf{不需要} teacher 速度场、不需要 $\Tmap$ 可逆或可微。
代价：丢失了 teacher 速度场在中间 $t$ 上的"软"信息，退化为"teacher 当数据生成器"。
\end{tcolorbox}

\eqref{eq:l0-cross} 与现有 \eqref{eq:l0-shared} 的差别：现有版本用 $\vT(z_t,t)$ 作
软靶子（teacher 在每个噪声水平上的意见）；跨 VAE 版本用解析 FM 靶子
$(\epsilon-z_0^S)$。当 $\Tmap$ 像素一致时，两者在 $\sigma_t\to0$ 极限一致；中间 $t$
的差异是"有损性"的来源之一。若坚持要软靶子，则需 \S\ref{sec:methods} 的线性传输
(M2) 经 Prop.~\ref{prop:affine}(b) 把 $\vT$ 直接以 $A\vT$ 搬运。

\subsection{L1：必须传输速度场，约束远强于 L0}

L1 要在 student 自己访问的状态 $x_k^S\in\Zs$ 上得到 teacher 的转移均值。由
Cor.~\ref{cor:l1-hard}，需要 $(\Tmap_\push\vT)(x_k^S)$，即可逆 + 可微 $\Tmap$。
三种可行度：

\begin{itemize}[leftmargin=1.4em]
\item \textbf{线性/仿射 $\Tmap=Az+b$（最干净）}：由 Prop.~\ref{prop:affine}(b)，
$\mu_T^S(x_k^S)=A\,\mu_T\big(A^{+}(x_k^S-b)\big)+b$（转移均值同样以仿射搬运）。
$A$ 行满秩可逆时无偏；$d_S\neq d_T$ 时用伪逆，引入投影损失。
\item \textbf{非线性 $\Tmap$（MLP，可行但贵）}：需要 $\Tmap^{-1}$（额外学一个逆网络
或用可逆架构）与 JVP（autograd 的 \texttt{jvp}）来算 $\Jac_{\Tmap}\vT$。逐时间步
代价 $\times$；且 Rem.~\ref{rem:nonlinear-break} 的路径破坏使闭式高斯 KL 的"同方差"
假设近似成立的条件变差。
\item \textbf{像素桥接 $\Tmap=\Enc_S\circ\Dec_T$（L1 最差）}：$\Tmap^{-1}$ 无解析式，
Jacobian 需穿过两个 VAE 的 VJP，逐步 decode+encode 代价 + 有损 + 不可逆，
\textbf{实践上不适合 L1}。
\end{itemize}

\begin{tcolorbox}[colback=green!4,colframe=green!50!black,title=关键洞察 2：L0/L1 对传输的要求是\emph{分层}的]
\textbf{L0 只需样本传输}（任何像素一致的 $\Tmap$ 都行，包括有损的像素桥接）；
\textbf{L1 需要速度场传输}（实际上要求仿射 $\Tmap$ 才干净）。因此一个自然的跨 VAE
XOPD 方案是：\emph{用便宜的像素桥接跑 L0 warmup，用学习式线性传输跑 L1}，两阶段用
不同的 $\Tmap$。这直接复用了 XOPD 已有的"L0 warmup $\to$ L1"两阶段开关。
\end{tcolorbox}

%% ============================================================
\section{算法族：六种传输的形式化与对比}
\label{sec:methods}
%% ============================================================

\subsection{M1：Pixel-space 桥接（decode-encode）}

\begin{equation}
\Tmap_{\mathrm{px}}=\Enc_S\circ\Dec_T,
\qquad
z_0^S=\Enc_S\big(\Dec_T(z_0^T)\big).
\end{equation}
\textbf{性质}：无需训练；像素一致性最好（$\delta_{\Tmap}$ 仅由两个 VAE 各自的重建
误差贡献）。\textbf{用于 L0}：直接（关键洞察 1）。\textbf{用于 L1}：差
（Cor.~\ref{cor:l1-hard}，不可逆 + 逐步 decode/encode 昂贵）。\textbf{代价}：每个
teacher 样本一次 $\Dec_T$ + 一次 $\Enc_S$（L0 中每 batch 一次，可接受）。对应用户
方法 1，是有损但最简单的 baseline。

\subsection{M2：学习式 latent projector（线性 / 仿射 / 非线性）}

学一个 $\Tmap_\phi:\Zt\to\Zs$，目标是像素一致：
\begin{equation}
\min_\phi\;
\E_{z\sim p_T}\Big[\big\lVert \Dec_S(\Tmap_\phi(z))-\Dec_T(z)\big\rVert^2\Big]
\;\;\Big(=\delta_{\Tmap_\phi}^2\text{ 的经验版}\Big),
\label{eq:m2-train}
\end{equation}
或在配对数据上直接 latent-latent 回归
$\E_{x}\lVert \Tmap_\phi(\Enc_T(x))-\Enc_S(x)\rVert^2$（更便宜，无需逐步解码）。
\textbf{线性 $\Tmap_\phi(z)=Az+b$}：由 Prop.~\ref{prop:affine}，L0 \emph{和} L1 都
干净（速度 $A\vT$、转移均值仿射搬运）；参数极少（$d_S\!\times\!d_T$），可闭式最小二乘
初始化。\textbf{非线性 MLP}：$\delta$ 更小但 L1 需逆+JVP（\S\ref{sec:l0l1}）。
对应用户方法 2，是本文推荐的\emph{L1 主力}。

\subsection{M3：CV-VAE 式重建对齐（把对齐推到 VAE 训练阶段）}

参考 \href{https://arxiv.org/abs/2405.20279}{CV-VAE (Zhao et al., 2024)}：与其在
推理时传输 latent，不如\textbf{重训一个与 teacher VAE "隐空间兼容"的 student VAE}。
CV-VAE 冻结预训练（image）VAE $(\Enc_i,\Dec_i)$，训练新（video）VAE
$(\Enc_v,\Dec_v)$，加两个像素重建正则：
\begin{equation}
\loss_{\mathrm{reg}}^{\mathrm{dec}}=\big\lVert \psi(X)-\Dec_i(\Enc_v(X))\big\rVert^2,
\quad
\loss_{\mathrm{reg}}^{\mathrm{enc}}=\big\lVert X-\Dec_v(\Enc_i(\psi(X)))\big\rVert^2,
\quad
\loss=\loss_{\mathrm{AE}}+\lambda_1\loss_{\mathrm{reg}}^{\mathrm{dec}}
+\lambda_2\loss_{\mathrm{reg}}^{\mathrm{enc}},
\label{eq:cvvae}
\end{equation}
（$\psi$ 是处理时序压缩的像素映射；他们消融出 decoder-only $\lambda_1{=}1,\lambda_2{=}0$
最好）。\textbf{迁移到我们的设定}：把 teacher VAE 当"冻结的预训练 VAE"，重训/微调
student VAE 使
$\Dec_T(\Enc_S(x))\approx x$ 且 $\Dec_S(\Enc_T(x))\approx x$，则两个隐空间近似
\emph{可互换}（$\Tmap\approx \mathrm{Id}$ 在通道对齐后），从源头消除传输需求。
\textbf{代价}：要训一个 VAE 且会改变 student 的 decoder（影响最终采样质量）；属于
"一次性重投资换取后续零传输开销"。\textbf{与 M1/M2 的关系}：M3 是把 $\delta_{\Tmap}$
\emph{在 VAE 训练阶段}压到极小，使 M1/M2 的传输几乎无损。

\subsection{M4（新提出）：概率耦合 / 分布矩匹配}

M1--M3 都是\emph{点对点}（sample-level）传输。但 XOPD 的 L1 本质是分布匹配（高斯
转移 KL）。我们提出在\emph{分布层面}对齐：不要求 $\Tmap(z_0^T)$ 与某个特定 student
隐变量逐点相等，只要求\textbf{推前分布}匹配 student 的边际：
\begin{equation}
\Tmap^\star=\arg\min_{\Tmap}\;
W_2\big(\Tmap_\push p_T,\;p_S\big)
\quad\text{或}\quad
\min_\Tmap \KL\big(\Tmap_\push p_T\,\|\,p_S\big),
\label{eq:m4-ot}
\end{equation}
其中 $p_T=\Enc_{T\push}p_{\mathrm{data}}$、$p_S=\Enc_{S\push}p_{\mathrm{data}}$ 是两个
VAE 在同一图像分布上的隐变量边际。\textbf{重要更正}：扩散用 VAE 的隐变量边际\emph{并不}
服从高斯（Rem.~\ref{rem:not-gaussian}），故 \eqref{eq:m4-ot} 的解\emph{一般不是}仿射的，
我们\textbf{不}声称"仿射是最优传输"。但可以退一步只对齐前两阶矩：取仿射
$\Tmap(z)=Az+b$ 使推前的均值/协方差与 $p_S$ 匹配（$\Tmap_\push p_T$ 与 $p_S$ 同均值
同协方差），其解为\textbf{矩匹配映射 (moment-matching map)}，由 teacher$\to$student 方向
的 Bures 公式给出：
\begin{equation}
A=\Sigma_T^{-1/2}\big(\Sigma_T^{1/2}\Sigma_S\,\Sigma_T^{1/2}\big)^{1/2}\Sigma_T^{-1/2},
\qquad
b=\mu_S-A\,\mu_T,
\label{eq:gaussian-ot}
\end{equation}
（$\mu_\bullet,\Sigma_\bullet$ 为两边际的均值/协方差，离线估计；通道数不同时在公共子空间上
做。可验证 $A\Sigma_T A^\top=\Sigma_S$，即协方差精确对齐。）矩匹配\textbf{不需要}分布高斯
——任何分布都有均值和协方差；它只在两边际\emph{恰好}近高斯时才升级为全局
$W_2$-最优传输 \eqref{eq:m4-ot}。\textbf{贡献}：给出 M2 线性传输一个\emph{有原则、无需
训练}的初始化（精确对齐二阶结构），并把"线性够不够"的判据交还给 $\delta_{\Tmap}$ 实测
（Rem.~\ref{rem:affine-geometry}）——而非依赖一个对扩散 VAE 不成立的高斯假设。

\subsection{M5（新提出）：cycle-consistent 联合学习传输}

同时学正/逆传输 $\Tmap_\phi:\Zt\to\Zs$ 与 $\Tmap_\psi:\Zs\to\Zt$，加 cycle 一致性，
直接满足 Cor.~\ref{cor:l1-hard} 对"可逆"的要求（无需对非线性 $\Tmap$ 求解析逆）：
\begin{equation}
\loss_{\mathrm{cyc}}
=\underbrace{\E\lVert \Dec_S\Tmap_\phi(z^T)-\Dec_T(z^T)\rVert^2}_{\text{像素一致}}
+\beta\,\underbrace{\E\lVert \Tmap_\psi\Tmap_\phi(z^T)-z^T\rVert^2}_{\text{cycle}}
+\beta\,\E\lVert \Tmap_\phi\Tmap_\psi(z^S)-z^S\rVert^2 .
\label{eq:m5}
\end{equation}
$\Tmap_\psi$ 提供 L1 所需的 $\Tmap^{-1}$（查询 teacher），Jacobian 经 autograd JVP。
\textbf{代价/收益}：比 M2 线性更表达，比纯非线性 M2 多了可逆保证；适合 $\delta$ 难以
被线性压低（两个 VAE 隐空间几何差异大）的情形。

\subsection{M6（新提出）：shared-decoder 微调（最小侵入的 M3）}

M3 重训整个 VAE 代价大。M6 只\textbf{冻结 student 的 $\Enc_S$、微调一个轻量
"适配解码头"} 使 teacher 隐变量能被 student 解码器读出：学
$\Dec_{S}^{\mathrm{adapt}}$ 使 $\Dec_S^{\mathrm{adapt}}(\Tmap(z^T))\approx\Dec_T(z^T)$，
其余 student 组件不动。介于 M2（不碰 VAE）与 M3（重训 VAE）之间。

\subsection{横向对比}

\begin{table}[h]
\centering\small
\begin{tabular}{@{}lcccccc@{}}
\toprule
& \textbf{需训练} & \textbf{$\delta_{\Tmap}$（有损）} & \textbf{L0} & \textbf{L1} &
\textbf{可逆/可微} & \textbf{额外开销} \\
\midrule
M1 pixel 桥接        & 否       & 中（2×VAE 重建） & \ding{51} & \ding{55}\,(贵+不可逆) & 否 & 每 batch decode+encode \\
M2 线性 projector    & 是（小） & 低（矩匹配残差） & \ding{51} & \ding{51}\,(干净)      & 是（伪逆） & 一次矩阵乘 \\
M2 非线性 projector  & 是（中） & 最低             & \ding{51} & $\triangle$\,(逆+JVP)  & 需额外逆网 & JVP/步 \\
M3 CV-VAE 重训       & 是（大） & 训后极低         & \ding{51} & \ding{51}              & $\approx$Id & VAE 训练（一次性） \\
M4 矩匹配（初始化）  & 闭式     & 低（二阶对齐）   & \ding{51} & \ding{51}\,(仿射)      & 是 & 离线估计 $\mu,\Sigma$ \\
M5 cycle 双向        & 是（中） & 低               & \ding{51} & \ding{51}              & 显式可逆 & 两网+JVP \\
M6 shared-decoder    & 是（小） & 低               & \ding{51} & \ding{51}              & --- & 解码头微调 \\
M7 白化/AdaLN 仿射   & 闭式     & 低（对角二阶）   & \ding{51} & \ding{51}\,(干净)      & 解析可逆 & 逐通道 $\gamma,\beta$ \\
\bottomrule
\end{tabular}
\caption{六种 VAE 隐空间传输方法对 XOPD L0/L1 的适配性对比。\ding{51}=干净适配，
$\triangle$=可行但昂贵，\ding{55}=不适合。}
\end{table}

\begin{tcolorbox}[colback=yellow!6,colframe=orange!60!black,title=推荐路线]
\textbf{先做 M1（L0）+ M2-线性（L1），并用 M4 的矩匹配闭式
\eqref{eq:gaussian-ot} 初始化 M2 的 $A,b$。} 理由：(i) M1 几乎零实现成本，立刻让
L0 warmup 可跑通，验证"teacher 数据是否对 student 有用"这一最关键问题；(ii) M2-线性
是唯一让 L1 干净成立的最小方案（几何结论，不依赖分布假设，Rem.~\ref{rem:affine-geometry}），
且 M4 给了它无需训练的二阶矩对齐初始化；
(iii) M3/M5/M6 作为后续提升项，当 $\delta_{\Tmap}$ 被测出过大（线性传输不足以对齐两个
隐空间）时再投入。
\end{tcolorbox}

%% ============================================================
\section{传输的初始化与 AdaLN 式可逆调制}
\label{sec:init-adaln}
%% ============================================================

本节回应两个具体设计问题：(1) transport 权重应如何初始化——恒等映射是否合适？
(2) 把 ODE 转成 SDE 后能否用 AdaLN 式的"高斯到高斯"全局调制来构造可逆传输？两个问题
都指向同一个对象——\textbf{逐通道仿射（对角白化）传输}——但都需要先澄清一个分布前提。

\subsection{初始化：恒等不是好默认，矩匹配才是"中性"初始化}

\begin{proposition}[恒等初始化的失效模式]
\label{prop:identity-bad}
把线性传输 $\Tmap(z)=Az+b$ 初始化为恒等（$A=I,b=0$）在跨 VAE 下一般\textbf{不可用}：
\begin{enumerate}[leftmargin=1.6em,label=(\alph*)]
\item \textbf{维度不合法}：$C_T\neq C_S$ 时 $A=I$ 无定义（FLUX.2 与 SD3.5 通道数不同）；
\item \textbf{尺度错位}：即使 $C_T=C_S$，两个 VAE 的 latent 一阶/二阶矩不同（SD3.5 用
``scaling\_factor``/``shift\_factor``，FLUX.2 用 BatchNorm 标准化），$A=I$ 把 teacher
latent 原样塞入 student 空间 $\Rightarrow$ OOD $\Rightarrow$ $\delta_{\Tmap}$ 巨大。
\end{enumerate}
\end{proposition}

正确的"中性初始化"是\textbf{对角矩匹配（白化-再着色）}：对齐逐通道均值/标准差
\begin{equation}
A_{\mathrm{diag}}=\mathrm{diag}\!\Big(\tfrac{\sigma_S^{(c)}}{\sigma_T^{(c)}}\Big),
\qquad
b=\mu_S-A_{\mathrm{diag}}\,\mu_T,
\label{eq:diag-mm}
\end{equation}
其中 $\mu_\bullet^{(c)},\sigma_\bullet^{(c)}$ 是两个 latent 边际的逐通道矩（离线估计，
$C_T\neq C_S$ 时在前 $\min(C_T,C_S)$ 个通道上做、其余零填充）。

\begin{remark}[为什么矩匹配是"恒等的正确推广"]
\label{rem:mm-is-identity}
\eqref{eq:diag-mm} 满足：当两个空间\emph{恰好相同}（同均值同方差）时
$A_{\mathrm{diag}}=I,b=0$，自动退化为恒等。因此"该不该用恒等初始化"的精确答案是：
\textbf{恒等只在退化（同空间）情形正确；一般情形应该用矩匹配——它在该退化时自己变成
恒等}。这也正是 \S\ref{sec:methods} M4 闭式 \eqref{eq:gaussian-ot} 的对角特例：M4 的
全协方差 Bures 解对齐二阶\emph{全}结构，\eqref{eq:diag-mm} 只对齐对角，更便宜、更稳。
\end{remark}

\begin{remark}[可训练非线性传输的"恒等初始化" = zero-init 残差]
\label{rem:zero-init}
若 transport 含可训练的\emph{非线性}部分 $g_\phi$，"do no harm" 初始化的正解不是把
$g_\phi$ 设成恒等，而是 \textbf{zero-init 残差}：
$\Tmap(z)=A_{\mathrm{diag}}z+b+g_\phi(z)$ 且 $g_\phi$ 末层零初始化（AdaLN-Zero /
ControlNet zero-conv / LoRA-B 的思路）。则训练\emph{初始}时 $\Tmap$ 恰好等于良态的仿射
baseline \eqref{eq:diag-mm}，非线性修正从 0 开始安全注入。这是想法 1 在非线性情形的正确
落地，并自然衔接下一小节。
\end{remark}

\subsection{ODE\,$\to$\,SDE 与 AdaLN 式调制：纠正前提，保留机制}

\paragraph{前提纠正：SDE 的高斯性在转移核，不在边际。}
"把 ODE 转成 SDE（Flow-SDE / CPS），中间步为高斯，于是两个 latent 空间的 latent 都是
高斯"——这里要区分两个分布（与 Rem.~\ref{rem:not-gaussian} 同源）：
\begin{itemize}[leftmargin=1.5em]
\item \textbf{转移核} $p(z_{t+dt}\mid z_t)$：SDE 下\emph{是}高斯（加性高斯噪声）。XOPD 的
$D_k$ 闭式高斯 KL 用的正是它——这\emph{不需要}把 ODE 转成 SDE，pathwise $D_k$ 在两种
dynamics 下都良定义。
\item \textbf{边际} $p(z_t)$，其中 $z_t=(1-\sigma_t)z_0+\sigma_t\epsilon$：\emph{非}高斯
（非高斯的 $p(z_0)$ 与高斯卷积），仅 $\sigma_t\to1$ 时趋于高斯。
\end{itemize}
所以 \textbf{ODE\,$\to$\,SDE 并不会使 latent 边际变高斯}；"两个空间的 latent 都是高斯"
这一前提\textbf{不成立}，"在两个高斯之间做变换"的字面表述需要放弃。

\paragraph{机制是对的：AdaLN 调制 $\equiv$ 可逆逐通道仿射（白化传输）。}
尽管前提不成立，想法 2 的\emph{机制}是有价值的：AdaLN 的逐通道 scale+shift
$\Tmap(z)=\gamma\odot z+\beta$ 恰是 Prop.~\ref{prop:affine} 的对角仿射，且
\begin{equation}
\Tmap^{-1}(z')=(z'-\beta)\oslash\gamma
\label{eq:adaln-inverse}
\end{equation}
\textbf{解析可逆}（$\gamma>0$），无需像 M5 那样单独学一个逆网络——这正好提供
Cor.~\ref{cor:l1-hard} 要求的 L1 逆映射。换言之，\textbf{想法 2 独立地重新发现了"仿射
调制是天然的可逆传输"}，且其参数恰由 \eqref{eq:diag-mm} 的矩匹配给出闭式初始化。我们把
这个对象命名为 \textbf{M7：WhiteningTransport（对角 AdaLN 仿射）}，作为 M2-线性的
对角特例：参数量 $2\min(C_T,C_S)$、可逆、零病态、矩匹配即得，是最稳的中性默认。

\begin{proposition}[时间条件对仿射传输无增益]
\label{prop:no-time-cond}
VAE 的 encode/decode 是与扩散时间 $\sigma$ \emph{无关}的固定映射，故两个 latent 空间的
关系本身与 $\sigma$ 无关。由 Prop.~\ref{prop:affine}(a)，一个\emph{静态}仿射
$\Tmap(z)=Az+b$ 已把\emph{所有} $\sigma$ 的状态正确传输（插值路径整体被仿射搬运）。因此
给仿射调制再加\emph{时间条件} $\gamma(\sigma),\beta(\sigma)$（AdaLN$(t)$）对仿射情形
\textbf{没有收益}——时间条件只有当 transport 是\emph{非线性}（不同 $\sigma$ 的边际形状差
异需要不同的非线性校正）时才可能有用。
\end{proposition}

\begin{remark}[一个细致的 caveat：均值干净，噪声需正交]
\label{rem:orthogonal-noise}
仿射传输对\emph{转移均值}的搬运是干净的（$D_k$ 用的就是均值，Prop.~\ref{prop:affine}(b)），
所以 pathwise 蒸馏不受影响。但若还想让传输后的\emph{噪声}也匹配 student 的各向同性标准
噪声，需要 $A$ 为正交矩阵（保协方差为各向同性）；而矩匹配/白化的 $A$ 一般\emph{非}正交，
故传输后的路径分布只是\emph{近似} on-distribution。这\textbf{不}影响 mean-matching 的
$D_k$ 使用，但它（而非时间条件）才是引入更富表达传输（M5 cycle / 正交约束）的真正动机。
若将来上 SDE+REINFORCE，则噪声项进入 log-prob，正交性才变得重要。
\end{remark}

\begin{tcolorbox}[colback=blue!4,colframe=blue!50!black,title=两个想法的结论]
\textbf{想法 1}：恒等初始化一般不可用（维度/尺度，Prop.~\ref{prop:identity-bad}）；正解是
\emph{矩匹配}中性初始化 \eqref{eq:diag-mm}（同空间时自动退化为恒等，
Rem.~\ref{rem:mm-is-identity}），非线性情形用 zero-init 残差
（Rem.~\ref{rem:zero-init}）。\\
\textbf{想法 2}：ODE$\to$SDE 不使 latent 边际高斯（前提需放弃），但 AdaLN 式逐通道仿射
调制是\emph{天然可逆}的传输（\eqref{eq:adaln-inverse}）、由矩匹配闭式初始化——即新增的
\textbf{M7 WhiteningTransport}。时间条件对仿射无增益（Prop.~\ref{prop:no-time-cond}），
故无需真的转 SDE；静态对角仿射即可。
\end{tcolorbox}


%% ============================================================
\section{实现设计（不写生产代码，只做架构）}
\label{sec:impl}
%% ============================================================

\subsection{改动面}

\begin{enumerate}[leftmargin=1.5em]
\item \textbf{放松硬约束}：\texttt{flux2\_klein.py:1363} 的 \texttt{in\_channels}
断言改为"共享 VAE 时断言相等，跨 VAE 时要求配置了 \texttt{vae\_transport}"。
\item \textbf{teacher VAE / 解码}：跨 VAE 时需要 teacher 的 $\Dec_T$（M1）或 teacher
$\Enc_T$（M2 训练）。新增 \texttt{load\_teacher\_vae(...)}，与现有
\texttt{load\_teacher\_text\_encoder} 一样\emph{仅预处理期常驻、之后 offload}（M1 的
$\Dec_T$ 用于 L0 生成靶子，可在预处理把 $z_0^S$ 直接缓存）。
\item \textbf{抽象 \texttt{VAETransport}}：一个小接口
\begin{verbatim}
class VAETransport(Protocol):
    def transport_sample(z_T) -> z_S          # M1/M2: 样本传输 (L0)
    def pushforward_velocity(v_T, z_S) -> v_S  # M2 线性: A v_T (L1)
    def inverse(z_S) -> z_T                    # L1 查询 teacher 用
\end{verbatim}
实现：\texttt{PixelBridgeTransport}(M1)、\texttt{LinearTransport}(M2/M4)、
\texttt{MLPTransport}(M2/M5)。
\item \textbf{L0 插入点}（\texttt{trainer.py:\_l0\_epoch}）：teacher rollout 得
$z_0^T$ 后，插入 $z_0^S=\Tmap.\texttt{transport\_sample}(z_0^T)$，其余 L0 不变
（\eqref{eq:l0-cross}，回归解析靶子或传输后的 teacher 速度）。
\item \textbf{L1 插入点}（\texttt{trainer.py:\_teacher\_next\_latents\_mean}）：teacher
forward 前用 $\Tmap.\texttt{inverse}$ 把 student 状态映回 $\Zt$ 查询，再用
\texttt{pushforward\_velocity} 把均值/速度映回 $\Zs$。
\item \textbf{config}：\texttt{vae\_transport: \{type: pixel|linear|mlp|none,
init: gaussian\_ot|lstsq|random, train: bool\}}；\texttt{none} 即现有共享-VAE 行为，
完全向后兼容。
\end{enumerate}

\subsection{两阶段用不同传输（呼应关键洞察 2）}
L0 用 \texttt{PixelBridgeTransport}（M1，无损训练、像素一致最好）；L1 用
\texttt{LinearTransport}（M2，矩匹配 \eqref{eq:gaussian-ot} 初始化、可选继续训练）。两者
共享同一个 teacher，互不冲突，且都走 XOPD 既有的 \texttt{use\_teacher\_transformer} +
DDP-bypass + autocast-cache 关闭路径（见 \texttt{CLAUDE.md} 的不变量）。

\subsection{先验证什么}
按 \S\ref{sec:formal} 的 $\delta_{\Tmap}$：实现任何传输前，\textbf{先离线测量}
$\delta_{\mathrm{px}}$（M1 桥接的像素往返误差）与矩匹配线性传输 \eqref{eq:gaussian-ot}
的残差 $\delta_{\mathrm{linear}}$，作为"该任务是否可蒸馏 / 线性是否够用"的前置体检
——若 $\delta_{\mathrm{px}}$ 已经很大（两个 VAE 语义不兼容），任何下游 XOPD 都救不回来，
应优先 M3/M6 对齐 VAE；若 $\delta_{\mathrm{px}}$ 小但 $\delta_{\mathrm{linear}}$ 大
（隐空间可互通但非线性关联），则线性 M2 不足，升级到非线性 M2 / M5。

%% ============================================================
\section{结论}
%% ============================================================

跨 VAE 的 XOPD 归结为设计隐空间传输 $\Tmap:\Zt\to\Zs$。核心理论结论：
\begin{itemize}[leftmargin=1.4em]
\item L0 与 L1 对 $\Tmap$ 的要求\emph{分层}：L0 只需样本传输（任何像素一致 $\Tmap$），
L1 需速度场推前（Prop.~\ref{prop:pushforward}），实际要求\emph{仿射}传输才干净
（Prop.~\ref{prop:affine}）。
\item 偏好仿射传输是\emph{几何}结论（Rem.~\ref{rem:affine-geometry}），\textbf{不}依赖
latent 服从高斯——事实上扩散 VAE 的隐变量\emph{边际是非高斯}的
（Rem.~\ref{rem:not-gaussian}）。M4 的闭式 \eqref{eq:gaussian-ot} 是\emph{矩匹配}初始化
（对齐二阶结构），只在边际恰好近高斯时才升级为 $W_2$-最优传输。
\item 所有方法的有损性可用单一\emph{经验}标量 $\delta_{\Tmap}$（\eqref{eq:delta}）量化
比较；"线性够不够"由 $\delta_{\mathrm{linear}}$ 实测判定，而非高斯定理。
\item 工程上：M1(L0)+M2 线性(L1)+M4 矩匹配初始化 是最小可行且理论自洽的首选；M3/M5/M6
作为 $\delta$ 过大时的升级项。
\end{itemize}

\end{document}
